\documentclass[11pt,a4paper]{article}
\usepackage[T1]{fontenc}
\usepackage[utf8]{inputenc}
\usepackage[english]{babel}
\usepackage{geometry,xcolor,listings,hyperref,amsmath}
\geometry{margin=23mm}
\definecolor{navy}{RGB}{18,49,88}
\lstset{language=Python,basicstyle=\ttfamily\small,breaklines=true,breakatwhitespace=false,
  frame=single,backgroundcolor=\color{black!3},columns=fullflexible,
  showstringspaces=false,keepspaces=true}
\hypersetup{colorlinks=true,linkcolor=navy}
\title{Sine MLP: line-by-line explanation}
\author{Frédéric Magoulès}
\date{November 2026}
\begin{document}
\maketitle
\noindent This document explains every nonblank line of the supplied \texttt{mlp\_sine\_intro.py}, without changing the program. Line numbers match the source exactly. Blank lines are omitted. Validation takes place after every epoch; final inference reuses validation inputs, without a separate test set.
\section*{Detailed explanations}
\subsection*{Line 1}
\begin{lstlisting}
"""Introductory MLP: fit sin(2*pi*x) with PyTorch on CPU."""
\end{lstlisting}
Module docstring: states the goal of fitting sin(2$\pi$x) with PyTorch on the CPU.

\subsection*{Line 2}
\begin{lstlisting}
from pathlib import Path
\end{lstlisting}
Imports Path to construct an operating-system-independent output path.

\subsection*{Line 3}
\begin{lstlisting}
import numpy as np
\end{lstlisting}
Imports NumPy, used here to set its random seed.

\subsection*{Line 4}
\begin{lstlisting}
import torch
\end{lstlisting}
Imports PyTorch for tensor operations and automatic differentiation.

\subsection*{Line 5}
\begin{lstlisting}
from torch import nn
\end{lstlisting}
Imports PyTorch neural-network layers and functions under the alias nn.

\subsection*{Line 6}
\begin{lstlisting}
from torch.utils.data import DataLoader, TensorDataset
\end{lstlisting}
Imports tools for datasets and mini-batch iteration.

\subsection*{Line 7}
\begin{lstlisting}
import matplotlib
\end{lstlisting}
Imports Matplotlib for plotting.

\subsection*{Line 8}
\begin{lstlisting}
matplotlib.use("Agg")
\end{lstlisting}
Selects the noninteractive Agg rendering backend before importing pyplot; figures can be saved without a GUI.

\subsection*{Line 9}
\begin{lstlisting}
import matplotlib.pyplot as plt
\end{lstlisting}
Imports the pyplot plotting interface.

\subsection*{Line 11}
\begin{lstlisting}
torch.manual_seed(7)
\end{lstlisting}
Sets PyTorch's pseudorandom seed to 7 for more reproducible initialization and shuffling.

\subsection*{Line 12}
\begin{lstlisting}
np.random.seed(7)
\end{lstlisting}
Also sets NumPy's pseudorandom seed.

\subsection*{Line 14}
\begin{lstlisting}
# Training and validation inputs are distinct.
\end{lstlisting}
Comment: training and validation input locations are different.

\subsection*{Line 15}
\begin{lstlisting}
x_train = torch.linspace(0, 1, 101).unsqueeze(1)
\end{lstlisting}
Creates 101 evenly spaced inputs from 0 to 1; unsqueeze(1) produces a column tensor of shape (101,1).

\subsection*{Line 16}
\begin{lstlisting}
x_val = torch.linspace(0.005, 0.995, 100).unsqueeze(1)
\end{lstlisting}
Creates 100 validation inputs from 0.005 to 0.995, offset from the training inputs.

\subsection*{Line 17}
\begin{lstlisting}
y_train = torch.sin(2 * torch.pi * x_train)
\end{lstlisting}
Computes the target values sin(2$\pi$x) at training inputs.

\subsection*{Line 18}
\begin{lstlisting}
y_val = torch.sin(2 * torch.pi * x_val)
\end{lstlisting}
Computes the exact target values at validation inputs.

\subsection*{Line 19}
\begin{lstlisting}
loader = DataLoader(TensorDataset(x_train, y_train), batch_size=32, shuffle=True)
\end{lstlisting}
Pairs inputs with targets and creates shuffled mini-batches of at most 32 samples.

\subsection*{Line 21}
\begin{lstlisting}
# One input, one output, and two small hidden layers.
\end{lstlisting}
Comment: one scalar input, one scalar output, and two hidden layers.

\subsection*{Line 22}
\begin{lstlisting}
model = nn.Sequential(nn.Linear(1, 32), nn.Tanh(), nn.Linear(32, 32),
\end{lstlisting}
Begins a sequential model: affine layer 1$\to$32, tanh activation, then affine layer 32$\to$32.

\subsection*{Line 23}
\begin{lstlisting}
                      nn.Tanh(), nn.Linear(32, 1))
\end{lstlisting}
Adds a second tanh activation and the affine output layer 32$\to$1; closes the model definition.

\subsection*{Line 24}
\begin{lstlisting}
optimizer = torch.optim.Adam(model.parameters(), lr=1e-2)
\end{lstlisting}
Creates an Adam optimizer over the model parameters with learning rate 0.01.

\subsection*{Line 25}
\begin{lstlisting}
loss_fn = nn.MSELoss()
\end{lstlisting}
Selects the mean squared error (MSE) loss function.

\subsection*{Line 26}
\begin{lstlisting}
train_loss, val_loss = [], []
\end{lstlisting}
Initializes lists that will record training and validation losses after each epoch.

\subsection*{Line 28}
\begin{lstlisting}
# Training
\end{lstlisting}
Marks the start of the training phase.

\subsection*{Line 29}
\begin{lstlisting}
for epoch in range(200):
\end{lstlisting}
Runs 200 epochs; each epoch is one full pass over the training data.

\subsection*{Line 30}
\begin{lstlisting}
    model.train()
\end{lstlisting}
Switches the model to training mode; especially relevant to dropout and batch normalization, neither of which is used here.

\subsection*{Line 31}
\begin{lstlisting}
    for xb, yb in loader:
\end{lstlisting}
Iterates through DataLoader mini-batches; xb holds inputs and yb the corresponding targets.

\subsection*{Line 32}
\begin{lstlisting}
        prediction = model(xb)           # Forward: predict y from x.
\end{lstlisting}
Feeds xb through the MLP to obtain predictions: the forward pass.

\subsection*{Line 33}
\begin{lstlisting}
        loss = loss_fn(prediction, yb)   # Loss: compare to sin(2*pi*x).
\end{lstlisting}
Computes mean squared error for the current mini-batch.

\subsection*{Line 34}
\begin{lstlisting}
        optimizer.zero_grad()           # Clear gradients from previous batch.
\end{lstlisting}
Clears gradients accumulated from previous optimization steps.

\subsection*{Line 35}
\begin{lstlisting}
        loss.backward()                 # Backward: compute gradients.
\end{lstlisting}
Backpropagates to compute derivatives of the loss with respect to trainable parameters.

\subsection*{Line 36}
\begin{lstlisting}
        optimizer.step()                # Update the MLP parameters.
\end{lstlisting}
Updates weights and biases using Adam.

\subsection*{Line 37}
\begin{lstlisting}
    # Validation
\end{lstlisting}
Marks the validation stage performed after every epoch.

\subsection*{Line 38}
\begin{lstlisting}
    model.eval()
\end{lstlisting}
Switches the model to evaluation mode.

\subsection*{Line 39}
\begin{lstlisting}
    with torch.no_grad():
\end{lstlisting}
Disables gradient tracking, reducing overhead during evaluation.

\subsection*{Line 40}
\begin{lstlisting}
        train_loss.append(loss_fn(model(x_train), y_train).item())
\end{lstlisting}
Computes MSE over the full training set and stores it as a Python scalar.

\subsection*{Line 41}
\begin{lstlisting}
        val_loss.append(loss_fn(model(x_val), y_val).item())
\end{lstlisting}
Computes and stores MSE on the distinct validation inputs.

\subsection*{Line 42}
\begin{lstlisting}
    if (epoch + 1) % 40 == 0:
\end{lstlisting}
Checks whether the one-based epoch number is divisible by 40.

\subsection*{Line 43}
\begin{lstlisting}
        print(f"Epoch {epoch+1:3d}: train={train_loss[-1]:.6f}, val={val_loss[-1]:.6f}")
\end{lstlisting}
Periodically prints training and validation losses to six decimal places.

\subsection*{Line 45}
\begin{lstlisting}
# Inference
\end{lstlisting}
Marks final inference: the trained parameters are no longer updated.

\subsection*{Line 46}
\begin{lstlisting}
model.eval()
\end{lstlisting}
Switches the model to evaluation mode.

\subsection*{Line 47}
\begin{lstlisting}
with torch.no_grad():
\end{lstlisting}
Disables gradient tracking during inference.

\subsection*{Line 48}
\begin{lstlisting}
    prediction = model(x_val).numpy().ravel()
\end{lstlisting}
Predicts on x\_val, converts the result to a NumPy array, and flattens it to one dimension.

\subsection*{Line 50}
\begin{lstlisting}
fig, axes = plt.subplots(1, 2, figsize=(10, 3.8))
\end{lstlisting}
Creates a figure with two side-by-side axes, sized 10 by 3.8 inches.

\subsection*{Line 51}
\begin{lstlisting}
axes[0].semilogy(range(1, 201), train_loss, label="Training")
\end{lstlisting}
Plots the training-loss history on a logarithmic vertical axis.

\subsection*{Line 52}
\begin{lstlisting}
axes[0].semilogy(range(1, 201), val_loss, label="Validation")
\end{lstlisting}
Plots the validation-loss history on the same axes.

\subsection*{Line 53}
\begin{lstlisting}
axes[0].set(xlabel="Epoch", ylabel="Mean squared error", title="Training history")
\end{lstlisting}
Sets axis labels and the first plot title.

\subsection*{Line 54}
\begin{lstlisting}
axes[0].legend()
\end{lstlisting}
Displays the first plot legend.

\subsection*{Line 55}
\begin{lstlisting}
axes[1].plot(x_val.numpy().ravel(), y_val.numpy().ravel(), label="Exact sine")
\end{lstlisting}
Plots the exact sine values at validation inputs.

\subsection*{Line 56}
\begin{lstlisting}
axes[1].plot(x_val.numpy().ravel(), prediction, "--", label="MLP prediction")
\end{lstlisting}
Plots the MLP predictions as a dashed line at those same inputs.

\subsection*{Line 57}
\begin{lstlisting}
axes[1].set(xlabel="x", ylabel="y", title="Function approximation")
\end{lstlisting}
Sets axis labels and the second plot title.

\subsection*{Line 58}
\begin{lstlisting}
axes[1].legend()
\end{lstlisting}
Displays the second plot legend.

\subsection*{Line 59}
\begin{lstlisting}
fig.tight_layout()
\end{lstlisting}
Automatically adjusts spacing to avoid overlapping plot elements.

\subsection*{Line 60}
\begin{lstlisting}
output = Path(__file__).with_name("mlp_sine_results.png") if "__file__" in globals() else Path("mlp_sine_results.png")
\end{lstlisting}
Chooses the output path: beside the script when \_\_file\_\_ exists, otherwise in the current directory (common in notebooks).

\subsection*{Line 61}
\begin{lstlisting}
fig.savefig(output, dpi=160)
\end{lstlisting}
Saves the PNG figure at 160 dots per inch.

\subsection*{Line 62}
\begin{lstlisting}
print(f"Saved {output}")
\end{lstlisting}
Prints the saved image path.

\subsection*{Line 63}
\begin{lstlisting}
plt.show()
\end{lstlisting}
Displays the figure when supported by the plotting environment.

\end{document}
